\chapter{已经验证到什么程度：数据、指标与下一步}
\label{ch:validation}

\section{验证必须先有可靠的物理参考}
神经网络预测快、训练误差小，不足以证明它在真实物理上正确。代理结果必须与完整 BFZI V150 电磁热直接法在\textbf{同一几何、相同材料、相同电压程序和相同时间点}比较。

物理参考自身也需要通过空间、时间与守恒检查。例如，分别使用两套海水 FV 网格，以及 1.875\,s 与 1\,s 两档真实内部热步，计算四条完整 600\,s 时程，再比较它们的动态阻抗与温度是否一致。被冻结的工况、阈值与采样时刻应保持不变。

早期出现过异常高温的工况，例如旧 F14 记录曾达到约 3441\,K。后续查明需要修正铜—封装热阻的几何表达，并检查单相海水与固态铜物性有效域。过高温度时程不能作为当前有效单相材料模型的训练参考。因此，项目先完成新的独立物理验收，再推进神经代理的广域重新评价。

\section{为什么要用动态阻抗增量误差}
对于一个给定工况，冷态阻抗往往已经占据完整 $Z$ 的大部分数值。如果按完整阻抗本身做相对误差，温度反馈预测得很差也可能得到一个很小的比例。

因此，我们关心的量是
\[
\Delta\mathbf z^*(t)=\mathbf z^*(t)-\mathbf z^*(0),
\]
其中 $\mathbf z$ 是 $Z_{11},Z_{12},Z_{22}$ 的六实分量。对跨网格比较和神经预测评价，分别使用
\begin{equation}
e_Z=
\frac{\left[\sum_t\|
\widehat{\Delta\mathbf z}(t)-\Delta\mathbf z^*(t)
\|_2^2\right]^{1/2}}
{\left[\sum_t\|\Delta\mathbf z^*(t)\|_2^2\right]^{1/2}}.
\end{equation}
直接法空间对照时，两套网格各自使用本身的冷态阻抗构造温变增量。报告还可以单独列出海水 $\Delta\zsea$ 误差，检查物理温度切线的实际质量。

温度采用整条时程、全部可比节点的最大绝对差（开尔文）。三维场采用相同空间坐标下的 $\mathbf J/q$ 相对 L2 或热点偏差；电路与热收支另行检查。

\section{S5 五阶近源积分物理直接法的最新冻结验收}
截至本文编写所依据的最新归档进度，12 个全新冻结工况 U01--U12 已完成两级网格、真实 1.875\,s 热步的空间对照，且全部满足原空间门槛。已完成的 41/48 条完整 600\,s 物理时程中，U01/U02/U03/U04/U05/U06/U07/U10 共八个工况完成全部四组合验收，其余四个工况待 1\,s 热步对照。

\begin{table}[htbp]\centering\small
\caption{S5 已完成空间验收的最大误差与原冻结门槛}
\begin{tabularx}{\textwidth}{Xrr}
\toprule
指标 & 目前最大值 & 原冻结门槛\\
\midrule
动态复阻抗跨网格相对 L2 & 1.4887\% & 2\%\\
铜温度最大跨网格差 & 0.6233\,K & 1\,K\\
双封装最大跨网格差 & 0.8916\,K & 1.5\,K\\
\bottomrule
\end{tabularx}
\end{table}

物理时间收敛的原阈值为动态阻抗 $\le1\%$、铜温度 $\le0.5$\,K、封装 $\le0.5$\,K，另需满足材料范围、端口功率和热能量守恒。U04 的五阶空间动态阻抗误差约为 1.032\%，U10 约为 0.838\%，两者相比此前三阶源积分明显改善。

\textbf{本书的 S5 数值是编写时的阶段快照，最终 12/12 状态以仓库的 \code{runs/bfzi\_physical\_unseen12\_s5.json} 和最新进度文件为准。} 未来通过后，还需要将新 S5 教师与神经模型的训练及独立测试保持一致。

\section{B 阶段：五材料热 ROM 曾通过什么验收}
在已锁定的历史 P03/P07 物理参考上，32 阶五材料热 ROM 分别将 2402 维热态降至 32 维。在独立电压和断电程序中，铜/封装最大温度附加误差很小，物理热能量预算通过对应阈值。

这说明三维物理温度能够由少量空间热模式有效表示。该阶段全程仍在每个热步调用原 BFZI 电磁计算，因此只替换热求解并未使整条时程大幅加速。

\section{C 阶段：物理切线残差和 SOE 记忆}
C 阶段在 P03/P07 上使用原 BFZI 冷态方向温度灵敏度、非线性铜阻抗曲线、32 阶热 ROM 和因果 SOE 残差网络，完成了新的 H1/H2 激励留出测试。

\begin{table}[htbp]\centering\small
\caption{C 阶段 H2 独立测试：各几何按 H1 选择的网络}
\begin{tabularx}{\textwidth}{lrrrr}
\toprule
几何 & 动态 $\Delta Z$ 误差 & 封装最大误差 & 铜最大误差 & 在线加速\\
\midrule
P03（分数阶） & 0.01531\% & 0.00907\,K & 0.00530\,K & 1135 倍\\
P07（单指数） & 0.04639\% & 0.07866\,K & 0.04517\,K & 1317 倍\\
\bottomrule
\end{tabularx}
\end{table}
上述加速指同台设备上该历史物理直接法与\textbf{在线时程推理}的时间比，未把首次冷态几何、热算子及方向切线构建时间算入在线时间。C 阶段具体使用较小的残差 \code{FracNet}，与第~\ref{ch:neural} 章的 40 维几何条件网络是两条不同实现路线。

\section{40 维几何条件网络和八模态高精度路线}
历史的纯神经 \code{GeomFrac} 在 P03/P07 整几何留出上得到约 3.47\% 与 3.87\% 的动态阻抗误差。结合原 BFZI 三维热 ROM 与八模态温度切线后，在 P07 的 150\,V、15\,s、完整 600\,s 独立历史中，完整动态复阻抗误差约为 0.0019\%--0.0026\%，温度与材料能量预算也得到相应的独立验证。

这些结果说明：\textbf{神经记忆提供跨历史预测能力，物理切线可以显著提升完整复阻抗精度}。同时，八模态校正需要离线方向求解，因此完整部署成本还应与跨几何共享路线比较。

\section{D 阶段：跨几何共享的成本与误差}
在 G1/G2 两个供体几何上构建热 ROM 和物理灵敏度后，对 P03、P07 和 G3 三个目标几何只进行一次低成本原 BFZI 冷态计算，使用共享热 ROM 与温度切线。其历史实测结果如下：
\begin{table}[htbp]\centering\small
\caption{D 阶段整几何留出及新几何部署成本}
\begin{tabularx}{\textwidth}{lrrr}
\toprule
目标 & 动态 $\Delta Z$ 误差 & 封装最大误差 & 新几何离线降本\\
\midrule
P03 & 1.3093\% & 0.1415\,K & 60.60\%\\
P07 & 0.9173\% & 0.6421\,K & 57.92\%\\
G3 & 0.8258\% & 0.3403\,K & 63.12\%\\
\bottomrule
\end{tabularx}
\end{table}
这些数字针对供体库已经建立之后的\textbf{每个新几何边际离线时间}。供体库的首次物理构建成本另有记录。当前共享路线仍需要目标几何的一次原 BFZI 冷态求解。

\section{为什么历史 24 组最终验收需要重新进行}
2026 年 11 月的早期广域 24 组验收中，旧代理面对较大的几何/材料/驱动变化时只通过其中 1 组，还有某些物理参考超过材料有效温度范围。该结果说明：早期模型的跨广域泛化尚未闭合，同时教师本身也需要更可信的数值与材料范围。

当前研究先解决物理教师的完整正确性，S5 冻结新样本覆盖了新的几何、激励及逐工况热参数，并保留原阈值。完成 12/12 后，再利用这一物理版本构建新教师，训练与重新评估跨几何共享、记忆网络和三维场预测，形成同一物理基准下的最终端到端验收。

\section{后续工作按什么顺序推进}
\begin{enumerate}
\item 完成 S5 的全部冻结样本，确认两级网格、两档时间步、材料范围和电热守恒均符合原标准；
\item 对需要高精度的神经参考更新冷态场、热 ROM、温度切线及教师时程，使训练和测试使用一致的物理求积方案；
\item 扩大训练中真实线圈形状、旋转、封装尺寸、热参数和启停历史的覆盖面，并保留整几何和整条历史的独立留出；
\item 在相同新物理参考上比较纯神经、八模态物理校正、共享 ROM 与残差网络，记录动态复阻抗、温度、三维 $J/q$、离线和在线成本；
\item 完成 E 阶段端到端系统验收，提供稳定的训练、保存、加载、推理和空间查询命令。
\end{enumerate}

本章对应仓库中的 \code{BFZI\_MAINLINE\_CLOSURE\_V1.md}、\code{PROJECT\_PROGRESS.md} 和 \code{evidence/} 下的 C、D、S5 阶段报告。
